\documentclass[11pt]{article}
\usepackage[a4paper,margin=1in]{geometry}
\usepackage[T1]{fontenc}
\usepackage{mathptmx}
\usepackage{amsmath,amssymb,amsthm,booktabs,enumitem,xcolor,fancyhdr,tabularx,array,graphicx}
\usepackage[ruled,vlined,linesnumbered]{algorithm2e}
\usepackage[section,above]{placeins}
\usepackage[hidelinks]{hyperref}
\usepackage{natbib}
\setcitestyle{numbers,square}

\pagestyle{fancy}\fancyhf{}
\renewcommand{\headrulewidth}{0.4pt}
\fancyhead[R]{\small Bluffing Against Strangers \quad \textsc{Multi-Agent Systems: Project Report}}
\fancyfoot[C]{\thepage}

\newtheorem{definition}{Definition}
\newtheorem{lemma}{Lemma}
\newtheorem{proposition}{Proposition}
\newtheorem{corollary}{Corollary}
\newtheorem{problem}{Problem}
\newcommand{\eps}{\varepsilon}
\newcommand{\epsbar}{\bar{\eps}}
\newcommand{\emax}{\eps_{\max}}
\newcommand{\BR}{\mathrm{BR}}
\makeatletter\newcommand{\tabin}[1]{\@@input #1 }\makeatother
% Data files live in ../results and ../figures when compiled inside the repository
% (report/final_report.tex), or in results/ and figures/ next to this file (e.g. Overleaf).
\IfFileExists{../results/numbers.tex}{\def\proot{../}}{\def\proot{}}
\input{\proot results/numbers.tex}% every result number in the text is generated by `python plots.py'
% ---- repository link ----
\newcommand{\repourl}{\url{https://github.com/Sejal-Kadgi/FAI}}
\setlist{nosep,leftmargin=*}
\setlength{\textfloatsep}{8pt plus 2pt minus 2pt}
\renewcommand{\topfraction}{0.9}\renewcommand{\bottomfraction}{0.8}\renewcommand{\textfraction}{0.07}
\renewcommand{\floatpagefraction}{0.8}\setcounter{topnumber}{3}\setcounter{bottomnumber}{2}\setcounter{totalnumber}{5}
\setlength{\floatsep}{6pt plus 2pt minus 2pt}

\title{\textsc{Bluffing Against Strangers: Learning Opponent Strategies\\ in a Short Encounter Under a Safety Constraint}\\[3pt]
\large Multi-Agent Systems: Project Report}
\author{Somireddy Bhavitha \and Sejal Vinayak Kadgi \and Siri Durugapu \and Lavanya Gupta\\
\small Indian Institute of Technology Gandhinagar \quad (Group 35)\\
\small \texttt{24110350@iitgn.ac.in, \{sejal.kadgi, siri.durugapu, lavanya.gupta\}@iitgn.ac.in}\\
\small Code and results: \repourl}
\date{November 2026}

\begin{document}
\maketitle
\thispagestyle{fancy}
\vspace{-1.5em}
\begin{abstract}
Equilibrium play in a hidden-information game is safe but leaves value unclaimed against opponents who bluff too much or call too rarely; exploiting them makes an agent exploitable. We study this trade-off in reduced two-player Bluff with undealt cards, where information about an opponent's honesty comes mainly from showdowns that cost points. We formulate \emph{safe exploitation in a short encounter}: over $G$ games, maximise the gain over a CFR+ equilibrium subject to a cap $\epsbar$ on the exploitability of \emph{each seat} in every game. The agent keeps an exact Bayesian posterior over a 25-type grid of bluffing/calling propensities and mixes a posterior best response with the equilibrium in realization-plan space. Convexity gives a cheap \emph{chord rule} that enforces the cap and certifies a per-seat floor on the expected payoff against any opponent. In exact experiments on trees of up to $\RnodesL$ nodes, the capped mixture never violated the cap and kept \RsafeShare\% of the oracle's gain (\RseatShare\% with per-seat weights), uncapped Bayesian learners beat bandit portfolios, and a restricted Nash response to the MAP type kept \RrnrShare\% with an observed (not enforced) worst-seat exploitability of at most $\RrnrEpsMaxE$.
\end{abstract}

\section{Introduction}
Whenever people negotiate, play poker, or defend a network, they face the same dilemma: a Nash equilibrium guarantees a minimum outcome whatever the opponent does~\cite{shoham2008}, but does nothing to profit from the opponent's habits, while a strategy tuned to one opponent can be punished by another~\cite{ganzfried2015}. The tension is sharpest in deception games. In Bluff (Cheat), players discard cards face down while claiming a required rank; only a call reveals whether a claim was honest, and a wrong call costs the caller the pile. Meeting a stranger for a short encounter, an agent has two scarce resources, the number of games and paid-for showdowns, and must decide how far to deviate from equilibrium given what it has learned, without exceeding a fixed exploitability budget.

\textbf{Related work.} CFR and CFR+~\cite{zinkevich2007,lanctot2009,neller2013} underlie superhuman poker agents~\cite{brown2018,brown2019}. Bayesian opponent modelling in poker goes back to~\cite{southey2005}; restricted Nash and data-biased responses trade exploitation for robustness~\cite{johanson2007,johanson2009}; safe exploitation is characterised in~\cite{ganzfried2015}; portfolio selection with bandits is studied in~\cite{bard2013,auer2002ucb,auer2002exp3}; see~\cite{albrecht2018} for a survey.

\textbf{Research questions.} RQ1: how fast are bluffing and calling propensities identified? RQ2: what is the best exploitation/exploitability trade-off at a given information level? RQ3: does Bayesian modelling beat model-free selection in a short encounter, and what does the safety cap cost? RQ4: do the conclusions survive opponents outside the model class?

\textbf{Changes after the idea review.} (1) Cards are left undealt, and E1 measures why this matters. (2) The deployed learner never stops learning; the budget is the encounter length, and freezing the belief survives only as a measurement device. (3) The exploitability limit is enforced exactly, in each seat separately (Proposition~1). (4) The scope is E1--E4 plus E5; type switching (E6) is future work.

\section{Model and Problem Formulation}
\begin{definition}[Reduced Bluff, Bluff$(K,m,h,T)$]
$K$ ranks with $m$ copies each; each player is dealt $h<Km/2$ cards and $U=Km-2h$ cards stay undealt and are never shown. At play $t$ player $t\bmod2$ puts $1\le c\le\min(m,|H|)$ cards face down claiming rank $r_t=t\bmod K$. The responder accepts or calls; on a call the cards are shown and the liar (or, if the claim was true, the caller) takes the pile. A player who empties their hand wins ($+1$) if the last play is accepted or called and true; after $T$ plays the player with fewer cards wins.
\end{definition}
\textbf{Visibility.} Cards are tracked by rank. A player sees its own hand and, for an opponent's play, only the claim size. A call shows the played cards to both players; the player who collects the pile sees its contents, the other only its size. Undealt cards and final hands are never shown. We keep two players and few cards because everything below is computed \emph{exactly} on the full tree. Testbeds: \textbf{S} $=$ Bluff(3,2,2,6) ($\RnodesS$ nodes) for all simulations and \textbf{L} $=$ Bluff(4,2,3,4) ($\RnodesL$ nodes) for the exact frontier experiment; both leave two cards undealt.

\begin{definition}[Opponent type]
$\theta=(\beta,\gamma)$: when both truthful and false plays exist, bluff w.p.\ $\beta$; call provably false claims always and other claims of size $c$ w.p.\ $\gamma^{1/c}$. Grid $\Theta=\{0.1,0.3,\dots,0.9\}^2$; archetypes Honest $(.1,.1)$, Aggressive bluffer $(.9,.1)$, Suspicious caller $(.3,.9)$, Random $(.5,.5)$, Paranoid liar $(.9,.9)$. The type is fixed for the whole match (an assumption; E5 tests opponents outside $\Theta$).
\end{definition}
\textbf{Exploitability.} The agent alternates seats, so its strategy is a pair $\pi=(\pi_1,\pi_2)$. Seat $i$'s exploitability is $\eps_i(\pi_i)=v_i^*-\min_{\pi_{-i}}u_i(\pi_i,\pi_{-i})$, and the safety metric is the worst seat, $\emax(\pi)=\max_i\eps_i(\pi_i)$; the seat average $\eps_{\mathrm{avg}}$ is used only for CFR convergence, since an average can hide one seat above the cap.

\begin{problem}[Safe exploitation in a short encounter]
Against a fixed unknown $\sigma_{\theta^*}$, over games $g=1..G$ with $\pi_g$ chosen from past observations,
$\max\ \mathbb E\big[\sum_g u(\pi_g,\sigma_{\theta^*})-u(\pi^*,\sigma_{\theta^*})\big]$ s.t.\ $\emax(\pi_g)\le\epsbar$ for every $g$.
\end{problem}
\textbf{Information level (measurement device).} Because the number of reveals depends on the agent's own play, E2--E3 also use beliefs $b^{(n)}$ recorded by a passive observer at its $n$-th reveal (a call, a claim provably false from its own hand, or a collected pile holding opponent cards). This freezing happens only in the analysis; deployed learners use every observation.

\section{Methodology}
\subsection{Equilibrium, best responses and exploitability}
The tree is built once and \emph{flattened}: every node is a row of numpy arrays and every edge out of a decision node is a global sequence index, so a two-seat behavioural strategy is one vector and every computation is a level-by-level vectorised pass. The reference $x^*$ is computed with \textbf{CFR+}: per infoset, regret matching $\sigma^t(a)\propto[R^t(a)]^+$, regrets updated as $R^{t+1}=\max(R^t+r^t,0)$ with counterfactual regrets $r^t$, and the output is the linearly weighted average strategy. A \textbf{best response} to any opponent reach weights $q(h)$ is one backward pass: at the agent's infosets pick $\arg\max_a\sum_{h\in I}\mathrm{CV}(ha)$, elsewhere sum children. \textbf{Exploitability} is two such passes.

\textbf{Exact vs approximate.} $x^*$ is only an approximate equilibrium ($\eps_{\mathrm{avg}}(x^*)=\RepsS$ on S, $\RepsL$ on L), so $v^*$ is not known exactly. Best responses, payoffs and exploitability of a given strategy are exact full-tree computations (agreement with the original implementation to $10^{-12}$). Since $v_1^*\in[\min_y u_1(x_1^*,y),\ \max_x u_1(x,x_2^*)]$, a bracket of width $2\eps_{\mathrm{avg}}(x^*)$ ($\RbracketS$ on S), we replace $v_i^*$ by the unfavourable end and obtain \emph{certified upper bounds} on each $\eps_i$; all caps and reported exploitabilities use them.

\subsection{Bayesian filter and posterior best response}
Own action probabilities are identical for all histories consistent with the agent's observations (perfect recall), so they cancel and the exact posterior with hidden cards marginalised is
\begin{equation}
b_{g+1}(\theta)\propto b_g(\theta)\sum_{h\in\mathcal H(O_g)}\sigma_c(h)\,\sigma_\theta(h).
\end{equation}
We precompute, per seat and type, $\sigma_c\sigma_\theta$ at every node and sum it over the agent's \emph{observation classes} (nodes it cannot tell apart); an update is a table lookup, and engine and filter share the same visibility rules by construction. The posterior best response $x^{\BR}_b$ is the best response to the reach weights $\sum_\theta b(\theta)\sigma_c\sigma_\theta(h)$, i.e.\ to a type drawn from $b$ at a hidden chance root.

\subsection{Safe mixing and the safety cap}
The agent plays $x_\alpha=(1-\alpha)x^*+\alpha x^{\BR}_b$ in realization-plan space~\cite{koller1996}.
\begin{lemma}[Mixture]
For a fixed opponent, the payoff is linear in $\alpha$, so the regret splits exactly into a \emph{price of safety} and a \emph{price of ignorance}: $r=(1-\alpha)G^*+\alpha\Delta$ with $G^*=u(\pi^{\BR}_{\theta^*})-u(\pi^*)$ and $\Delta=u(\pi^{\BR}_{\theta^*})-u(\pi^{\BR}_b)\ge0$. Each $\eps_i(x_{\alpha,i})$ is convex in $\alpha$ (a maximum of linear functions), hence so is $\emax(x_\alpha)$.
\end{lemma}
\begin{proposition}[Exact enforcement]
If $\eps^*=\emax(x^*)\le\epsbar$, the feasible weights form an interval $[0,\alpha_{\max}]$ and the \emph{chord rule}
\[\hat\alpha=\min\{1,\ (\epsbar-\eps^*)/(\emax(x^{\BR})-\eps^*)\}\]
is feasible; capping any weight rule at $\hat\alpha$ keeps both seats within $\epsbar$ in every game.
\end{proposition}
\begin{proof}A sublevel set of a convex function on $[0,1]$ containing 0 is an interval, and $\emax(x_{\hat\alpha})\le(1-\hat\alpha)\eps^*+\hat\alpha\,\emax(x^{\BR})=\epsbar$.\end{proof}
\begin{corollary}[Per-seat safety floor]
If $\emax(\pi_g)\le\epsbar$ then $u_i(\pi_{g,i},\sigma)\ge v_i^*-\epsbar$ in each seat and against \emph{every} opponent $\sigma$, since $u_i(\pi_i,\sigma)\ge\min_{\pi_{-i}}u_i=v_i^*-\eps_i(\pi_i)$.
\end{corollary}
The guarantee bounds the \emph{expected} payoff of each game, not individual outcomes; it is relative to $v_i^*$, not to what $x^*$ would earn against that opponent; and it needs no correct opponent model. Three cap methods are implemented: the common chord rule above; a \emph{per-seat chord} (seat constraints decouple, so $\hat\alpha_i$ uses $\eps_i$ only and $\hat\alpha_i\ge\hat\alpha$); and \emph{bisection} on $\emax$ (tolerance $1/64$). If $\epsbar<\emax(x^*)$ the code raises an error instead of silently playing $\alpha=0$. Algorithm~\ref{alg} summarises the learner.

\begin{algorithm}[t]
\caption{Safe Bayesian exploitation in an encounter of $G$ games}\label{alg}
\small
\KwIn{$x^*$, $\eps^*=\emax(x^*)$, grid $\Theta$, prior $b_0$, $\epsbar$, weight rule $\alpha^{\text{rule}}$, $\lambda$}
\lIf{$\epsbar<\eps^*$}{\textbf{report infeasible} and stop}
$b\gets b_0$\;
\For{$g=1,\dots,G$}{
  $x^{\BR}\gets$ best response to $\sum_\theta b(\theta)\sigma_c\sigma_\theta$, plus $\lambda\,\mathrm{EIG}_b$ on call edges\;
  $\hat\alpha\gets$ chord rule$(\epsbar,\eps^*,\emax(x^{\BR}))$ \tcp*{Prop.~1; per seat: $\hat\alpha_1,\hat\alpha_2$}
  $\alpha\gets\min\{\alpha^{\text{rule}}(b),\hat\alpha\}$;\ $x\gets(1-\alpha)x^*+\alpha x^{\BR}$\tcp*{$\eps_1,\eps_2\le\epsbar$}
  play game $g$ in seat $1+(g\bmod2)$ with $x$; update $b$ by Eq.~(1) with the whole game\;
}
\end{algorithm}

\textbf{Information-seeking calls.} At each agent call edge the best response gets a bonus $\lambda\,\mathrm{EIG}_b(I)$, the exact expected reduction in type entropy from the showdown; it shapes $x^{\BR}$ before mixing, so Proposition~1 still applies. $\lambda=0.2$ was fixed in advance and not tuned (the bonus is at most $0.2\ln25\approx0.64$ payoff units).

\subsection{Learners compared and protocol}
\textbf{Reference}: Equilibrium $x^*$; Oracle (best response to the true opponent; defines regret). \textbf{Bayesian, unconstrained}: Bayes-greedy ($\alpha=1$), Bayes-MAP (best response to the MAP type), Bayes-gated ($\alpha=1-H(b)/H(b_0)$). \textbf{Bayesian, safety-capped}: Bayes-safe ($\alpha=\hat\alpha$), Bayes-safe-seat (per-seat chord), Bayes-safe-gated, Bayes-safe-active ($\lambda=0.2$). \textbf{Model-free}: UCB1, Exp3~\cite{auer2002ucb,auer2002exp3} and sliding-window $\epsilon$-greedy over the portfolio $\{x^*\}\cup\{\BR(\sigma_\theta):\theta\text{ an archetype}\}$. \textbf{Robust game-theoretic}: RNR-MAP, a restricted Nash response~\cite{johanson2007} ($p=0.5$, CFR+ in the game where the opponent is forced to the model w.p.\ $p$) to the current MAP type, precomputed for all 25 types. Only the safety-capped group enforces $\epsbar$.

\textbf{Protocol.} $\epsbar=0.05$, $G=100$ games (E4/E5), 30 seeds, uniform prior, paired deals across learners (separate random stream). For every game we record the \emph{exact} expected payoff of the committed strategy, its per-seat exploitability, $\alpha$ and reveals. Each pass is $O(N)$; a CFR+ iteration takes $\approx\RmsIterS$\,ms on S and $\approx\RmsIterL$\,ms on L; the filter stores $O(|\Theta|N)$ numbers per seat. The whole pipeline (\texttt{python run\_all.py}) took \RpipelineHours{} hours on two CPU cores and reproduced all results bit for bit from scratch.

\section{Results}
\textbf{E1: equilibrium and undealt cards.} CFR+ converges far faster than CFR (Fig.~\ref{fig:e1}; $\RepsS$ on S after \RitersS{} iterations). Table~\ref{tab:fb1} confirms the reviewers' point: under a full deal \RfullProvT\% (T=4) and \RfullProvSix\% (T=6) of equilibrium lies are provably false to the responder; with two cards undealt lies make up \RpartLieLo--\RpartLieHi\% of claims and only \RpartProvLo--\RpartProvHi\% are provable, so bluffing becomes a real strategic choice.

\begin{figure}[!htbp]\centering\includegraphics[width=\linewidth]{\proot figures/e1_cfr.pdf}
\vspace{-1.2em}\caption{E1: seat-averaged exploitability of the CFR and CFR+ average strategies.}\label{fig:e1}\end{figure}
\begin{table}[!htbp]\centering\footnotesize
\resizebox{0.9\linewidth}{!}{\begin{tabular}{@{}llrrrrrrr@{}}\toprule
Game & deal & undealt & opp.\ hands & lies (\%) & lies provable (\%) & decidable (\%) & calls (\%) & $v_1^*$\\\midrule
\tabin{\proot results/e1_fb1.tex}
\bottomrule\end{tabular}}
\caption{E1 (feedback 1), under each game's equilibrium: share of lies the responder can prove false, of claims decidable without calling, and opponent hands consistent with one's own.}\label{tab:fb1}\end{table}

\textbf{E2: identification (RQ1).} The call propensity is learned about three times faster than the bluff propensity (Fig.~\ref{fig:e2}): after 10 games the posterior-mean error is $\RgamErrTen$ for $\gamma$ versus $\RbetErrTen$ for $\beta$, and $\RgamErrFifty$ versus $\RbetErrFifty$ after 50, because $\gamma$ shows in every response while $\beta$ needs a reveal. Payoff-equivalence speeds things up: for three archetypes \ReqSetSize{} of 25 types have near-optimal best responses ($\eta=0.01$), while Aggressive bluffer needs about \RtauAggr{} games. Information-seeking calls barely changed the number of reveals (\RrevPas{} vs \RrevAct{} per game).

\begin{figure}[!htbp]\centering\includegraphics[width=\linewidth]{\proot figures/e2_ident.pdf}
\vspace{-1.2em}\caption{E2: posterior mass on the payoff-equivalent set (solid passive, dashed information-seeking) and error of the posterior means of $\beta$ and $\gamma$ against games and against reveals.}\label{fig:e2}\end{figure}

\textbf{E3: exploitation--exploitability frontier (RQ2).} Linearity of the gain holds to $\RlinErrS$ and each seat's exploitability is convex in $\alpha$; on S it is even numerically linear (second differences $\le\RcurvS$), so the chord rule attains the tight $\alpha_{\max}$ (obtained by interpolating $\emax$ over 21 grid values of $\alpha$), while on L it keeps \RchordTightL\% of it. At low information the price of ignorance dominates the regret, at high information the price of safety (Fig.~\ref{fig:voi}). Under a fixed cap, more information can therefore \emph{lower} the safe gain on S (Fig.~\ref{fig:e3}): regret at $\epsbar=0.05$ grows from \RregZero{} per game at $n=0$ to \RregInf{} at $n=\infty$, because a vague belief gives a hedged, low-exploitability best response ($\emax\approx\RepsPrior$) whereas a sharp one is extreme in at least one seat (up to \RepsWorstS). On L, where best responses are exploitable in both seats, the effect reverses (Honest: \RhonZeroL{} to \RhonInfL; Fig.~\ref{fig:e3l}). Re-solving beats mixing: with perfect knowledge the mixture captures \RmixHon\%, \RmixAgg\% and \RmixPar\% of $G^*$ for Honest, Aggressive bluffer and Paranoid liar (\RseatHon\%, \RseatAgg\%, \RseatPar\% per seat), RNR \RrnrHon\%, \RrnrAgg\% and \RrnrPar\% at worst-seat exploitability $\le\RrnrEpsMaxS$; on L the mixture keeps \RmixLlo--\RmixLhi\%, RNR \RrnrLlo--\RrnrLhi\%.

\begin{figure}[!htbp]\centering\includegraphics[width=\linewidth]{\proot figures/e3_frontier_small.pdf}
\vspace{-1.2em}\caption{E3 on S: best safe gain $F_n(\epsbar)$ of the mixture for beliefs built from $n$ reveals (common weight; dashed: per-seat weights at $n=\infty$) and the restricted Nash response with the true model (dotted); dashed grey: $G^*$.}\label{fig:e3}\end{figure}

\begin{figure}[!htbp]\centering
\includegraphics[width=0.82\linewidth]{\proot figures/e3_voi_small.pdf}
\vspace{-0.8em}\caption{E3 on S. Left: unconstrained gain of the posterior best response relative to $G^*$ as information grows. Right: Lemma~1's price of safety and price of ignorance at $\epsbar=0.05$ (largest feasible common $\alpha$, mean over archetypes).}\label{fig:voi}\end{figure}
\begin{figure}[!htbp]\centering\includegraphics[width=\linewidth]{\proot figures/e3_frontier_large.pdf}
\vspace{-1.2em}\caption{E3 on the larger testbed L $=$ Bluff(4,2,3,4), same layout as Fig.~\ref{fig:e3}.}\label{fig:e3l}\end{figure}

\textbf{E4: learner comparison (RQ3).} Table~\ref{tab:e4} groups the learners. \emph{Unconstrained}: Bayes-greedy keeps \RgreedyShare\% of the oracle's gain (regret \RgreedyRegret{} over 100 games) and beats UCB1 by $\RgreedyVsUCB$ per game; the model-free learners keep \RmfLo--\RmfHi\% and lose to equilibrium against two opponents. None limits exploitability (worst seat up to \RgreedyEpsMax). \emph{Constrained}: the capped learners never exceeded $\epsbar$ in either seat in \RcappedGames{} games (maximum $\RcappedMax$). Bayes-safe keeps \RsafeShare\% (equal to unconstrained UCB1, $\RsafeVsUCB$), Bayes-safe-seat \RseatShare\% ($+\RseatVsSafe$), because best responses are asymmetric across seats (Honest: $\eps_1=\RhonSeatOne$, $\eps_2=\RhonSeatTwo$). \emph{RNR-MAP} keeps \RrnrShare\% with worst-seat exploitability $\le\RrnrEpsMaxE$ (Fig.~\ref{fig:scatter}); this is an observation, not a guarantee, as nothing in RNR ties its exploitability to $\epsbar$, but it shows that re-solving against the model exploits mistakes in parts of the tree an equilibrium opponent never reaches. Gating and information seeking did not pay ($-\RactiveVsSafe$ per game for $\lambda=0.2$).

\begin{table}[!htbp]\centering\footnotesize\setlength{\tabcolsep}{3.5pt}
\begin{tabular}{@{}lrrrrr|rrrr@{}}\toprule
& \multicolumn{5}{c|}{mean gain over $\pi^*$ per game} & & & \multicolumn{2}{c}{worst-seat $\emax$}\\
Learner & Honest & Aggr. & Susp. & Random & Paran. & mean & $R_{100}$ & mean & max\\\midrule
\tabin{\proot results/e4_table.tex}
\bottomrule\end{tabular}
\caption{E4 on S ($G=100$, 30 seeds; 95\% CIs of per-opponent means $\le\RciMaxE$). Only the safety-capped group enforces $\emax\le\epsbar=0.05$.}\label{tab:e4}\end{table}

\textbf{E5: misspecification (RQ4).} Against off-grid types (Table~\ref{tab:e5}), Bayes-greedy and RNR-MAP stay close to the oracle (regret $\le\RoffRegMax$ over 100 games vs \RoffEqLo--\RoffEqHi{} for equilibrium play). Structurally different opponents (hand-size-dependent, greedy) remove most of the gain (\RsizeGreedy{} of \RsizeOracle{} for Bayes-greedy), but not the safety: in every game and seat the capped learners' expected payoff stayed within $\RfloorCapped$ of $v_i^*$ (guarantee: 0.05), while unconstrained UCB1 fell $\RfloorUCB$ below and lost $\RucbVsEq$ per game to the equilibrium opponent. Per-seat capping again beat the common weight (\RseatShareE\% vs \RsafeShareE\% of the oracle gain).

\begin{table}[!htbp]\centering\footnotesize\setlength{\tabcolsep}{3pt}
\resizebox{\linewidth}{!}{\begin{tabular}{@{}lrrrrrrr|rrrr@{}}\toprule
& \multicolumn{7}{c|}{mean gain over $\pi^*$ per game} & & & \multicolumn{2}{c}{worst-seat $\emax$}\\
Learner & (.2,.8) & (.6,.2) & (.8,.6) & Size-dep. & Max-dump & Noisy eq. & Equilibrium & mean & $R_{100}$ & mean & max\\\midrule
\tabin{\proot results/e5_table.tex}
\bottomrule\end{tabular}}
\caption{E5 on S: off-grid types $(\beta,\gamma)$; a player whose bluffing depends on hand size and calling on claim size (Size-dep.); a greedy largest-claim player (Max-dump); $0.7\pi^*+0.3$ uniform (Noisy eq.); and the equilibrium itself.}\label{tab:e5}\end{table}
\begin{figure}[!htbp]\centering
\includegraphics[width=0.49\linewidth]{\proot figures/e4_scatter.pdf}\hfill
\includegraphics[width=0.49\linewidth]{\proot figures/e5_scatter.pdf}
\vspace{-0.8em}\caption{Mean gain against mean worst-seat exploitability of the strategies played, E4 (left) and E5 (right); dashed: $\epsbar$.}\label{fig:scatter}\end{figure}

\textbf{Discussion and limitations.} H2 and H4 are supported, and H1 in a stronger form (exploitability linear on S). H3 is partly refuted: the decomposition behaves as predicted, but on S the safe gain falls with information. H5 holds: zero per-seat violations, at a cost that per-seat weights largely recover, while re-solving (RNR) dominates mixing. Limits: exact computation restricts us to reduced two-player games; types are two-parameter and fixed; $\epsbar$, $\lambda$ and $p$ were not tuned; exploitability is a worst case that scripted opponents never realise.

\section{Contribution}
\begin{center}\small
\begin{tabularx}{\linewidth}{@{}l X l@{}}\toprule
Member & Assigned tasks & Status\\\midrule
Somireddy Bhavitha & Game engine with undealt cards and flattened tree; CFR/CFR+; exact best response and exploitability; E1. & Fulfilled\\
Sejal Vinayak Kadgi & Opponent library; exact Bayesian filter over observation classes; posterior/MAP best response; E2. & Fulfilled\\
Siri Durugapu & Realization-plan mixing; chord, per-seat and bisection caps; information-seeking bonus; RNR; E3. & Fulfilled\\
Lavanya Gupta & Bandit, no-model and RNR-MAP baselines; paired-seed runner, tests and plotting; E4 and E5. & Fulfilled\\
All members & Problem formulation, analysis, report. & Fulfilled\\\bottomrule
\end{tabularx}\end{center}

\begin{thebibliography}{99}\small\setlength{\itemsep}{0pt}
\bibitem{albrecht2018} S.~V. Albrecht and P.~Stone. Autonomous agents modelling other agents: A comprehensive survey and open problems. \emph{Artificial Intelligence}, 258:66--95, 2018.
\bibitem{auer2002ucb} P.~Auer, N.~Cesa-Bianchi, and P.~Fischer. Finite-time analysis of the multiarmed bandit problem. \emph{Machine Learning}, 47:235--256, 2002.
\bibitem{auer2002exp3} P.~Auer, N.~Cesa-Bianchi, Y.~Freund, and R.~E. Schapire. The nonstochastic multiarmed bandit problem. \emph{SIAM J.\ Computing}, 32(1):48--77, 2002.
\bibitem{bard2013} N.~Bard, M.~Johanson, N.~Burch, and M.~Bowling. Online implicit agent modelling. In \emph{AAMAS}, 2013.
\bibitem{brown2018} N.~Brown and T.~Sandholm. Superhuman AI for heads-up no-limit poker: Libratus beats top professionals. \emph{Science}, 359:418--424, 2018.
\bibitem{brown2019} N.~Brown and T.~Sandholm. Superhuman AI for multiplayer poker. \emph{Science}, 365:885--890, 2019.
\bibitem{ganzfried2015} S.~Ganzfried and T.~Sandholm. Safe opponent exploitation. \emph{ACM TEAC}, 3(2):8, 2015.
\bibitem{johanson2009} M.~Johanson and M.~Bowling. Data biased robust counter strategies. In \emph{AISTATS}, 2009.
\bibitem{johanson2007} M.~Johanson, M.~Zinkevich, and M.~Bowling. Computing robust counter-strategies. In \emph{NeurIPS}, 2007.
\bibitem{koller1996} D.~Koller, N.~Megiddo, and B.~von Stengel. Efficient computation of equilibria for extensive two-person games. \emph{Games and Economic Behavior}, 14(2):247--259, 1996.
\bibitem{lanctot2009} M.~Lanctot, K.~Waugh, M.~Zinkevich, and M.~Bowling. Monte Carlo sampling for regret minimization in extensive games. In \emph{NeurIPS}, 2009.
\bibitem{neller2013} T.~W. Neller and M.~Lanctot. An introduction to counterfactual regret minimization. 2013.
\bibitem{shoham2008} Y.~Shoham and K.~Leyton-Brown. \emph{Multiagent Systems}. Cambridge University Press, 2008.
\bibitem{southey2005} F.~Southey et al. Bayes' bluff: Opponent modelling in poker. In \emph{UAI}, 2005.
\bibitem{zinkevich2007} M.~Zinkevich, M.~Johanson, M.~Bowling, and C.~Piccione. Regret minimization in games with incomplete information. In \emph{NeurIPS}, 2007.
\end{thebibliography}
\end{document}